\chapter{Ambiente e Simulazione}

\textsc{E.L.B.E.R.R.}\ (\textit{Efficient Logistics by Exploration with Robotic Retrieval}) \`e una
simulazione multi-agente di recupero e consegna oggetti in un magazzino a pianta variabile.
Cinque agenti autonomi operano in parallelo su una griglia 2D con visione e comunicazione
locali, senza alcun coordinamento centralizzato.

\section{Griglia}

L'ambiente \`e rappresentato da una griglia quadrata $N \times N$ le cui celle appartengono a
uno dei seguenti tipi:

\begin{itemize}
  \item \texttt{EMPTY} (0) --- corridoio libero, percorribile dagli agenti;
  \item \texttt{WALL} (1) --- ostacolo o scaffale, non percorribile;
  \item \texttt{WAREHOUSE} (2) --- cella interna al magazzino (zona di deposito/ricarica);
  \item \texttt{ENTRANCE} (3) --- porta d'ingresso del magazzino;
  \item \texttt{EXIT} (4) --- porta d'uscita del magazzino.
\end{itemize}

Gli agenti possono attraversare solo celle \texttt{EMPTY}, \texttt{ENTRANCE} e \texttt{EXIT};
le porte dei magazzini sono state gestite in modo che siano percorribili esclusivamente come destinazione finale di un percorso, mai come
transito intermedio. 
Questo perchè è stato osservato che in alcuni casi, gli agenti consideravano erroneamente le porte come scorciatoie dovute alla presenza di troppi agenti nelle vicinanze, causando ulteriore congestione e collisioni evitabili. 
\section{Magazzini}

Ogni magazzino \`e una zona rettangolare con un ingresso e un'uscita separati. L'istanza viene
caricata da un file JSON che specifica la pianta, le posizioni dei magazzini e la
distribuzione casuale degli oggetti. Gli agenti consegnano gli oggetti entrando
dall'ingresso, depositandoli su una cella interna e uscendo dall'uscita.

\section{Ciclo di simulazione}

Ad ogni tick il simulatore esegue, in ordine:

\begin{enumerate}
  \item \textbf{Sense} --- ogni agente aggiorna la propria mappa locale tramite ray casting
        nel raggio di visione $v_r$;
  \item \textbf{Communicate} --- gli agenti entro il raggio di comunicazione reciproco
        $c_r$ si scambiano mappe locali, oggetti noti e posizioni degli agenti conosciuti
        (merge bidirezionale);
  \item \textbf{Decide} --- ogni agente sceglie la prossima mossa in base alla propria
        strategia di esplorazione;
  \item \textbf{Act} --- le mosse vengono applicate con gestione delle collisioni;
  \item \textbf{Pickup/Deliver} --- gli agenti raccolgono oggetti e li consegnano;
  \item \textbf{Metriche} --- aggiornamento dei contatori.
\end{enumerate}

La simulazione termina quando tutti gli oggetti sono stati consegnati oppure si raggiunge
il limite di 500-750 tick.

\section{Gestione delle collisioni}

La prevenzione delle collisioni è implementata con un approccio \emph{ibrido distribuito}:
ogni agente tenta attivamente di evitarle basandosi su informazioni locali, mentre il
simulatore funge da arbitro imparziale per i conflitti inevitabili.

\subsection*{Autogestione locale}

Ad ogni tick, nel costruire la propria destinazione, ogni agente filtra le celle
attraverso la funzione di pathfinding utilizzando l'insieme delle posizioni occupate
(\texttt{known\_agents}) entro una finestra di validit\`a temporale di 5~tick.
Agenti che comunicano direttamente si scambiano inoltre posizioni note, propagando
informazioni sullo stato del sistema via \textit{gossip protocol}.

Questo meccanismo permette ai singoli agenti di \emph{evitare proattivamente} collisioni
di cui sono a conoscenza, manifestando una forma di coordinamento localizzato e
biologicamente realistica.

\subsection*{Arbitrato centrale}

Tuttavia, in un sistema decentralizzato con comunicazione locale limitata, non è possibile
garantire che ogni agente conosca la posizione in tempo reale di tutti gli altri. Conflitti
emergono inevitabilmente quando:

\begin{itemize}
  \item Due agenti che non si vedono e non hanno comunicato tentano la stessa destinazione;
  \item L'informazione su un agente è scaduta (oltre 5~tick).
\end{itemize}

In questi casi, il simulatore applica una risoluzione deterministica \emph{prima di eseguire
alcun movimento}, garantendo che non si producano mai stati incoerenti:

\begin{enumerate}
  \item Un agente che non si muove (\texttt{dest\,=\,None}) \emph{blocca} la propria cella:
        nessun altro può entrarvi.
  \item Se due o più agenti richiedono la stessa cella, vince chi trasporta un oggetto
        (priorità funzionale); a parità di stato, vince l'ID minore (arbitrio deterministico
        per garantire no-deadlock).
        Gli altri restano fermi e potranno riprovare al prossimo tick.
  \item Le porte (\texttt{ENTRANCE} e \texttt{EXIT}) sono automaticamente aggiunte
        all'insieme protected durante la pianificazione, evitando che gli agenti le attraversino
        innecessariamente.
\end{enumerate}

\section{Prenotazione degli ingressi ai magazzini}

Quando un agente trasporta un oggetto, deve scegliere a quale ingresso consegnarlo.
La selezione non \`e statica: ad ogni tick viene rivalutata tramite un costo dinamico
(descritto nel Capitolo~3) che combina distanza, congestione locale e prenotazioni note.

Per evitare oscillazioni, la scelta viene vincolata con un \emph{lock di 8~tick}:
l'agente cambia destinazione solo se trova un ingresso con costo
almeno il 20\% inferiore a quello attuale.
La prenotazione viene condivisa via gossip durante la comunicazione, cos\`\i\
gli altri agenti possono tenerne conto nel calcolo della congestione.

\section{Comunicazione}

Due agenti comunicano se la distanza di Manhattan tra le loro posizioni \`e $\leq \min(c_{r,a}, c_{r,b})$.
Lo scambio \`e bidirezionale e include:

\begin{itemize}
  \item Mappa locale (merge delle celle esplorate);
  \item Insieme degli oggetti noti non ancora raccolti;
  \item Posizioni note degli altri agenti (si conserva la pi\`u recente);
  \item Prenotazioni di consegna ai magazzini (anti-congestione).
\end{itemize}

Non esiste un canale broadcast globale: la conoscenza si propaga solo per contatto diretto,
creando un effetto di \textit{gossip protocol} distribuito.
